\section{Validation and Results}
\section*{Validation introduction}
Validation was performed at several levels because a successful compilation
does not prove that the learning lifecycle is correct. The checks below cover
automated backend tests, frontend builds, API behavior, database consistency
and browser journeys for the three roles.

\section{Validation strategy}
Validation combined unit and integration tests for the backend, an Angular
production build, Swagger checks and live HTTP scenarios. The objective was
not only to verify that an endpoint returned a response, but to verify the
business rule represented by that response.

\begin{table}[H]\centering
\caption{Validation activities}\label{tab:validation}
\begin{tabularx}{\textwidth}{>{\bfseries}p{0.30\textwidth}X}
\toprule Activity & Result\\\midrule
Backend clean tests & Passed after CRUD, grading and deletion changes.\\
Angular production build & Passed with the final dashboard and branding changes.\\
OpenAPI availability & Swagger UI and OpenAPI JSON returned successfully.\\
Module deletion & Returned successful deletion after dependent cleanup.\\
Answer persistence & Review returned the selected option for the correct question.\\
Duplicate submission & A second attempt was rejected with HTTP 409.\\
User deletion & Dependent learner records were removed without a foreign-key failure.\\
\bottomrule
\end{tabularx}
\end{table}

\section{Validation traceability}
\begin{table}[H]\centering
\caption{Traceability between requirements and validation}
\label{tab:traceability}
\begin{tabularx}{\textwidth}{>{\bfseries}p{0.25\textwidth}X}
\toprule Requirement & Validation evidence\\\midrule
Security & JWT-aware login design, role-aware routes and protected operations.\\
Consistency & Real module and user deletion scenarios without foreign-key failure.\\
Traceability & Detailed answer review, submission date and approval/unapproval.\\
Maintainability & Successful layered backend tests and Angular production build.\\
AI continuity & Manual quiz creation remains available beside AI generation.\\
One-time assessment & First submission succeeds and repeated submission returns 409.\\
Material availability & Files remain accessible after quiz submission.\\
\bottomrule
\end{tabularx}
\end{table}

\section{Functional scenarios}
The first scenario creates a module and attaches documents. The second
generates or manually creates a quiz and publishes it. The learner then opens
the module, selects answers and submits once. The instructor sees the
submission in the review area, expands the detailed answers and approves it.
The administrator can consult the grouped submission statistics.

\section{Quality results}
The principal defects discovered during development were incorrect answer
keys, misleading HTTP 403 responses caused by error forwarding, incomplete
dependent deletion and an interface that allowed a second attempt. Each issue
was fixed at its root: question IDs are now used for answer maps, application
errors return explicit statuses, deletion is transactional, and both backend
and frontend prevent retakes.

\section{Security and operational limitations}
The JWT integration provides the intended authentication and authorization
boundary for deployment. Secret values, including the AI provider key, must
remain in environment configuration and must be rotated if exposed. File
uploads should receive size, type and content scanning limits before production.
Rate limiting and audit retention are also recommended for a public
deployment.

\section{Difficulties, limitations and perspectives}
The most important difficulties were not isolated coding errors but
interactions between features. Deleting a module required understanding all
dependent records. Persisting an answer required a stable question key rather
than a temporary array index. The approval workflow required a reversible
state, not only a one-way button. These issues were corrected through service
rules, repository operations and end-to-end checks.

The current validation does not constitute a production load test or an
independent security audit. Future work should measure performance with a
larger catalogue, add automated accessibility checks, scan uploaded files,
rotate refresh tokens, export audit records and add notifications. The AI
workflow should also display confidence or source references and be evaluated
against a larger curated question corpus.

\section{User experience observations}
The live screenshots show a coherent visual language: blue-gray surfaces,
strong contrast, shield branding, status badges and action-oriented cards.
The instructor screen emphasizes operational work, while the learner screen
keeps the path to course material and assessment simple. The administrator
screen prioritizes counts, grouped submissions and governance actions.

\section{Personal contribution}
This internship required work across requirements analysis, Angular
components, Spring services, relational integrity, third-party AI integration,
API documentation and testing. The most valuable lesson was that a feature is
not complete when its happy path works: its deletion behavior, error response,
role boundary and persistence format must also be verified.


